\chapter{Bredhjet e rastit, difuzioni dhe zinxhirët Markov}
\begin{qellime}
Studenti lidh bredhjen e rastit me difuzionin; nxjerr ligjin e zhvendosjes katrore mesatare; simulon bredhje në një dhe dy përmasa; përkufizon zinxhirin Markov, shpërndarjen stacionare, kthyeshmërinë dhe kohën e korrelacionit.
\end{qellime}

\section{Bredhja simetrike në një përmasë}
Le të jenë hapat $S_i=\pm\ell$ me probabilitet $1/2$. Pas $N$ hapash,
\begin{equation}
 X_N=\sum_{i=1}^{N}S_i,qquad \E[X_N]=0,qquad \E[X_N^2]=N\ell^2.
\end{equation}
Zhvendosja tipike rritet si $\sqrt N$, jo si $N$. Nëse çdo hap zgjat $\Delta t$, atëherë
\begin{equation}
 \avg{X^2}=2Dt,qquad D=\frac{\ell^2}{2\Delta t}.
\end{equation}
Në $d$ përmasa, $\avg{r^2}=2dDt$.

\begin{figure}[ht]
\centering\includegraphics[width=.95\textwidth]{21_bredhja_e_rastit.pdf}
\caption{Realizime të bredhjes dhe rritja lineare e zhvendosjes katrore mesatare.}
\end{figure}

\section{Kufiri i vazhdueshëm dhe difuzioni}
Probabiliteti diskret plotëson ekuacionin master
\begin{equation}
 p_j^{n+1}=\tfrac12p_{j-1}^n+\tfrac12p_{j+1}^n.
\end{equation}
Zhvillimi Taylor dhe kufiri $\ell,\Delta t\to0$ me $\ell^2/(2\Delta t)=D$ japin
\begin{equation}
 \frac{\partial p}{\partial t}=D\frac{\partial^2p}{\partial x^2}.
\end{equation}
Për një grimcë fillimisht në origjinë,
\begin{equation}
 p(x,t)=\frac1{\sqrt{4\pi Dt}}\exp\left(-\frac{x^2}{4Dt}\right).
\end{equation}

\begin{figure}[ht]
\centering\includegraphics[width=.74\textwidth]{22_difuzioni_gaussian.pdf}
\caption{Profili i difuzionit zgjerohet, ndërsa normalizimi ruhet.}
\end{figure}

\section{Kufij, drift dhe kohë e parë kalimi}
Një kufi reflektues ruan grimcën; një kufi thithës e heq kur arrihet. Koha e parë e kalimit është koha kur trajektorja prek për herë të parë një objektiv. Me drift $v$,
\begin{equation}
 \frac{\partial p}{\partial t}=-v\frac{\partial p}{\partial x}+D\frac{\partial^2p}{\partial x^2}.
\end{equation}
Kjo lidh transportin determinist me luhatjet. Në mjedise heterogjene, hapi ose probabiliteti mund të varen nga pozicioni.

\section{Zinxhirët Markov}
Një proces diskret është Markov nëse e ardhmja, me kusht gjendjen e tanishme, nuk varet nga historia:
\begin{equation}
 P(X_{n+1}=j|X_n=i,X_{n-1},\ldots)=P_{ij}.
\end{equation}
Matrica $P$ ka elemente jo-negative dhe rreshta me shumë një. Shpërndarja evoluon si $\vect p_{n+1}=\vect p_nP$. Shpërndarja stacionare plotëson
\begin{equation}
 \vect\pi=\vect\pi P.
\end{equation}
Një zinxhir i fundmë, i pareduktueshëm dhe aperiodik ka shpërndarje stacionare unike dhe konvergon drejt saj.

\begin{figure}[ht]
\centering\includegraphics[width=.74\textwidth]{23_zinxhiri_markov.pdf}
\caption{Humbja graduale e kujtesës së gjendjes fillestare në një zinxhir me dy gjendje.}
\end{figure}

\section{Kthyeshmëria dhe balanca e detajuar}
Kushti
\begin{equation}
 \pi_iP_{ij}=\pi_jP_{ji}
\end{equation}
quhet balancë e detajuar. Ajo garanton stacionaritetin, por nuk është kusht i domosdoshëm. Interpretimi është ekuilibër i fluksit probabilitar ndërmjet çdo çifti gjendjesh.

\section{Korrelacioni dhe madhësia efektive}
Në një zinxhir, mostrat pasuese janë të korreluara. Për një observabël $A$,
\begin{equation}
 \rho_A(k)=\frac{\Cov(A_n,A_{n+k})}{\Var(A)},\qquad
 \tau_{\mathrm{int}}=\frac12+\sum_{k=1}^{\infty}\rho_A(k).
\end{equation}
Varianca e mesatares sillet afërsisht si
\begin{equation}
 \Var(\bar A)\approx\frac{2\tau_{\mathrm{int}}}{N}\Var(A),
\end{equation}
pra $N_{\mathrm{eff}}\approx N/(2\tau_{\mathrm{int}})$. Rritja e numrit të hapave nuk ndihmon shumë kur zinxhiri përzihet ngadalë.

\begin{lidhje}
Laboratori simulon bredhje 1D dhe 2D, mat $\avg{r^2}$, zbaton kufij thithës e reflektues dhe vlerëson kohën e korrelacionit të një zinxhiri të vogël.
\end{lidhje}

\section*{Pyetje kontrolli}
\begin{enumerate}
\item Si del ekuacioni i difuzionit nga ekuacioni master?
\item Çfarë dallimi ka stacionariteti nga pavarësia?
\item Pse $N$ hapa Markov nuk japin $N$ mostra të pavarura?
\end{enumerate}
\section{Algoritmet e bredhjes dhe të zinxhirit Markov}
\subsection{Bredhja dydimensionale}
\begin{lstlisting}[language=Python,caption={Shumë realizime të bredhjes në 2D.}]
def bredhje_2d(rng, n_realizime, n_hapa):
    drejtimi = rng.integers(0, 4, size=(n_realizime, n_hapa))
    dx = (drejtimi == 0) - (drejtimi == 1)
    dy = (drejtimi == 2) - (drejtimi == 3)
    x = np.c_[np.zeros(n_realizime), np.cumsum(dx, axis=1)]
    y = np.c_[np.zeros(n_realizime), np.cumsum(dy, axis=1)]
    return x, y

rng = np.random.default_rng(2026)
x, y = bredhje_2d(rng, 50_000, 2_000)
msd = np.mean(x*x + y*y, axis=0)
\end{lstlisting}
Përshtatja e $\avg{r^2}$ kundrejt kohës jep $4D$ në dy përmasa. Pikat e hershme dhe efektet e kufirit trajtohen veçmas.

\subsection{Kufiri thithës dhe koha e parë e kalimit}
\begin{lstlisting}[language=Python]
def koha_e_pare(x, prag):
    prekje = np.abs(x) >= prag
    ka = prekje.any(axis=1)
    t = np.full(len(x), np.nan)
    t[ka] = np.argmax(prekje[ka], axis=1)
    return t
\end{lstlisting}
Vlerat pa prekje janë të censuruara dhe nuk duhet të zëvendësohen me kohën maksimale pa e deklaruar.

\subsection{Evolucioni me matricë kalimi}
\begin{lstlisting}[language=Python,caption={Shpërndarja stacionare e një zinxhiri të fundmë.}]
P = np.array([[0.90, 0.10], [0.25, 0.75]])
p = np.array([1.0, 0.0])
for _ in range(100):
    p = p @ P
print(p, p @ P - p)

w, V = np.linalg.eig(P.T)
i = np.argmin(np.abs(w - 1.0))
pi = np.real(V[:, i]); pi /= pi.sum()
\end{lstlisting}

\subsection{Autokorrelacioni dhe $N_{\mathrm{eff}}$}
\begin{lstlisting}[language=Python,caption={Vlerësim fillestar i kohës së integruar të autokorrelacionit.}]
def autocorr(x, max_lag):
    x = np.asarray(x) - np.mean(x)
    c = np.correlate(x, x, mode="full")[len(x)-1:]
    c = c[:max_lag+1] / c[0]
    return c

rho = autocorr(observabli, 500)
pozitive = rho[1:][rho[1:] > 0]
tau_int = 0.5 + pozitive.sum()
N_eff = len(observabli)/(2*tau_int)
\end{lstlisting}
Prerja në vonesën e parë negative është rregull i thjeshtë, jo universal. Për analiza serioze përdoret dritare automatike ose bllokim.

\begin{lidhje}
Notebook-u duhet të paraqesë njëkohësisht trajektoret, $\avg{r^2}$, shpërndarjen e pozicionit dhe kohën e parë të kalimit. Secila mat një aspekt të ndryshëm të procesit.
\end{lidhje}

Për proceset e rastit dhe terminologjinë e statistikës janë konsultuar \cite{capollari,butka,tasho}; trajtimi i zinxhirëve Markov dhe autokorrelacionit mbështetet te \cite{robert,krauth}.

\subsection{Bredhja e rastit dhe kufiri i difuzionit}
Në një përmasë, hapi $\xi_n$ merr vlerat $\pm\ell$ me probabilitete të barabarta dhe
\begin{equation}
 X_N=\sum_{n=1}^{N}\xi_n.
\end{equation}
Pavarësia jep $\E[X_N]=0$ dhe $\Var(X_N)=N\ell^2$. Nëse një hap zgjat $\Delta t$, atëherë $t=N\Delta t$ dhe
\begin{equation}
 \E[X(t)^2]=\frac{\ell^2}{\Delta t}t=2Dt,
 \qquad D=\frac{\ell^2}{2\Delta t}.
\end{equation}
Për të marrë kufirin e vazhdueshëm mbahen $D$ dhe $t$ të fiksuara ndërsa $\ell,\Delta t\to0$. Shpërndarja e pozicionit i afrohet
\begin{equation}
 p(x,t)=\frac{1}{\sqrt{4\pi Dt}}\exp\left(-\frac{x^2}{4Dt}\right),
\end{equation}
që zgjidh $p_t=Dp_{xx}$. Ky është një shembull themelor ku një model mikroskopik stokastik çon në ekuacion makroskopik determinist.

Me drift, probabilitetet nuk janë të barabarta. Në kufirin e duhur shfaqet ekuacioni adveksion--difuzion
\begin{equation}
 p_t=-v p_x+Dp_{xx}.
\end{equation}
Kushtet thithëse përfaqësojnë largimin e grimcës; kushtet reflektuese ruajnë probabilitetin. Koha e parë e kalimit është ndryshore rasti dhe trajektoret që nuk e arrijnë kufirin brenda kohës së simulimit janë të censuruara, jo zero.

\subsection{Zinxhirët Markov dhe shpërndarja stacionare}
Për një zinxhir me hapësirë të fundme dhe matricë kalimi $P$, shpërndarja rresht evoluon si
\begin{equation}
 \vect\pi_{n+1}^T=\vect\pi_n^TP.
\end{equation}
Shpërndarja stacionare plotëson $\vect\pi^TP=\vect\pi^T$, $\pi_i\ge0$ dhe $\sum_i\pi_i=1$. Nëse zinxhiri është i pakthyeshëm dhe aperiodik, shpërndarja konvergon te stacionarja pavarësisht gjendjes fillestare. Shpejtësia lidhet me boshllëkun spektral ndërmjet vlerës vetjake $1$ dhe vlerës së dytë në modul.

Balanca e detajuar,
\begin{equation}
 \pi_iP_{ij}=\pi_jP_{ji},
\end{equation}
është kusht i mjaftueshëm për stacionaritet, sepse mbledhja sipas $i$ jep $(\vect\pi^TP)_j=\pi_j$. Ai nuk është kusht i domosdoshëm: mund të ketë rryma stacionare që nuk anulohen çift më çift.

\subsection{Autokorrelacioni dhe madhësia efektive}
Për një seri stacionare $A_n$, autokovarianca është
\begin{equation}
 C(k)=\E[(A_n-\mu)(A_{n+k}-\mu)],
 \qquad \rho(k)=C(k)/C(0).
\end{equation}
Varianca e mesatares për $N$ vëzhgime është
\begin{align}
 \Var(\overline A)&=\frac{1}{N^2}\sum_{i,j}\Cov(A_i,A_j)\\
 &\simeq\frac{C(0)}{N}\left[1+2\sum_{k=1}^{\infty}\rho(k)\right]
 =\frac{2\tau_{\mathrm{int}}C(0)}{N},
\end{align}
ku
\begin{equation}
 \tau_{\mathrm{int}}=\frac12+\sum_{k=1}^{\infty}\rho(k).
\end{equation}
Prandaj
\begin{equation}
 N_{\mathrm{eff}}=\frac{N}{2\tau_{\mathrm{int}}}.
\end{equation}
Formula $s/\sqrt N$ nënvlerëson gabimin kur mostrat janë të korreluara. Shuma e autokorrelacionit duhet prerë në një vonesë ku zhurma e vlerësimit nuk dominon; alternativa praktike janë mesataret në blloqe dhe shumë zinxhirë të pavarur.

\subsection{Koha e përzierjes dhe diagnostika}
Koha e përzierjes mat sa shpejt shpërndarja e zinxhirit harron gjendjen fillestare. Heqja e një numri fiks hapash pa diagnostikë nuk garanton termalizim. Vëzhgohen trajektoret, mesataret kumulative dhe krahasimi i zinxhirëve nga gjendje të ndryshme. Hollimi i zinxhirit zakonisht hedh informacion; ruajtja e të gjitha mostrave dhe përdorimi i $N_{\mathrm{eff}}$ është më transparent, përveç kufizimeve të ruajtjes ose kostos së vlerësimit.

\subsection{Përmbledhje dhe ushtrime}
Bredhja e rastit lidh fluktuacionet mikroskopike me difuzionin. Zinxhiri Markov shton varësinë kohore dhe kërkon korrigjimin e gabimit standard me autokorrelacionin.
\begin{enumerate}
\item Nxirrni $\E[r^2]=2dDt$ për bredhjen në $d$ përmasa.
\item Zgjidhni kohën mesatare të daljes nga një interval me kufij thithës.
\item Gjeni shpërndarjen stacionare dhe vlerat vetjake të një matrice kalimi $3\times3$.
\item Simuloni një zinxhir me dy gjendje dhe krahasoni $N_{\mathrm{eff}}$ teorik me atë të vlerësuar.
\item Tregoni me shembull se stacionariteti nuk nënkupton balancë të detajuar.
\end{enumerate}

\subsection{Koha mesatare e daljes}
Për bredhjen simetrike në nyjet $0,1,\ldots,L$ me kufij thithës, le të jetë $T_i$ numri mesatar i hapave deri në dalje kur nisemi nga $i$. Kushtëzimi ndaj hapit të parë jep
\begin{equation}
 T_i=1+\frac12T_{i-1}+\frac12T_{i+1},
 \qquad T_0=T_L=0.
\end{equation}
Riorganizimi jep diferencën e dytë
\begin{equation}
 T_{i+1}-2T_i+T_{i-1}=-2.
\end{equation}
Zgjidhja kuadratike që plotëson kufijtë është
\begin{equation}
 T_i=i(L-i).
\end{equation}
Në kufirin e difuzionit, koha mesatare $\tau(x)$ plotëson ekuacionin e prapëm
\begin{equation}
 D\tau''(x)=-1,qquad \tau(0)=\tau(a)=0,
\end{equation}
me zgjidhje $\tau(x)=x(a-x)/(2D)$. Kjo lidh drejtpërdrejt ekuacionin diskret me modelin e vazhdueshëm.

Nëse një kufi është reflektues, kushti bëhet derivat zero. Me drift, operatori i prapëm është $v\partial_x+D\partial_{xx}$. Kohët e para të kalimit janë të ndjeshme ndaj bishtave dhe kërkojnë shumë realizime; mesatarja vetëm mund të mos përshkruajë shpërndarjen.

\subsubsection{Spektri i matricës së kalimit}
Për një zinxhir të kthyeshëm, matrica e kalimit është vetë-adjunkte në prodhimin skalar të peshuar
\begin{equation}
 \langle f,g\rangle_\pi=\sum_i\pi_if_ig_i.
\end{equation}
Vlerat vetjake janë reale dhe plotësojnë $1=\lambda_1>\lambda_2\ge\cdots\ge-1$. Komponentët jo stacionarë shuhen si $\lambda_k^n$. Boshllëku spektral $1-\max_{k>1}|\lambda_k|$ kontrollon relaksimin më të ngadaltë.

Për një observabël që projektohet fort mbi modin e dytë, autokorrelacioni mund të sillet si $\rho(n)\simeq\lambda_2^n$. Atëherë
\begin{equation}
 \tau_{\mathrm{int}}\simeq\frac12+\frac{\lambda_2}{1-\lambda_2}
 =\frac{1+\lambda_2}{2(1-\lambda_2)}.
\end{equation}
Kjo tregon lidhjen ndërmjet përzierjes, boshllëkut spektral dhe gabimit të mesatares.

\subsubsection{Vlerësimi praktik i autokorrelacionit}
Vlerësuesi i autokovariancës në vonesën $k$ është
\begin{equation}
 \widehat C(k)=\frac{1}{N-k}\sum_{n=1}^{N-k}
 (A_n-\overline A)(A_{n+k}-\overline A).
\end{equation}
Për $k$ të madh mbeten pak çifte dhe vlerësuesi bëhet i zhurmshëm. Shuma naive deri në $N$ mund të luhatet fort. Përdoret një dritare vetëkonsistente: shuma ndalet kur $k$ është disa herë më i madh se $\widehat\tau_{\mathrm{int}}$, ose kur autokorrelacioni hyn në një rajon të pajtueshëm me zhurmën.

Metoda e blloqeve ndan serinë në blloqe me gjatësi $B$ dhe llogarit mesataren e çdo blloku. Kur $B\gg\tau$, mesataret e blloqeve janë afërsisht të pavarura. Gabimi i vlerësuar duhet të arrijë një pllajë me rritjen e $B$, para se numri i blloqeve të bëhet tepër i vogël.

\subsubsection{Ergodiciteti dhe kurthet}
Stacionariteti nuk mjafton nëse zinxhiri nuk mund të arrijë të gjithë mbështetjen. Një zinxhir i reduktueshëm mund të ketë disa shpërndarje stacionare dhe rezultati varet nga gjendja fillestare. Periodiciteti mund të shkaktojë alternim pa konvergjencë pikë për pikë, edhe kur mesataret kohore sillen më mirë.

Në hapësira me moda të ndara nga barriera, zinxhiri mund të jetë teorikisht ergodik, por praktikisht të mos kalojë barrierën gjatë simulimit. Trajektorja duket e qëndrueshme brenda një mode dhe jep pacaktueshmëri artificialisht të vogël. Nisja e shumë zinxhirëve, propozimet globale, temperimi ose metodat e grupimeve përdoren për të zbuluar dhe zbutur këtë problem.

\subsubsection{Renditja e faqeve dhe zinxhiri Markov}
Në renditjen e faqeve, kalimi ndjek një lidhje me probabilitet $\alpha$ dhe ``teleportohet'' me probabilitet $1-\alpha$. Termi i teleportimit e bën matricën pozitive, kështu që teorema Perron--Frobenius garanton një vektor stacionar unik pozitiv. Parametri $\alpha$ balancon respektimin e strukturës së lidhjeve me shpejtësinë e përzierjes.

Faqet pa dalje duhet të trajtohen: kolona përkatëse zëvendësohet me një shpërndarje probabilitare. Pa këtë hap, matrica nuk është stohastike dhe masa humbet. Ky shembull tregon se pastrimi i të dhënave është pjesë e modelit matematik, jo vetëm detaj programimi.

